% Standalone explanatory derivation of the implemented construction.
% This document describes the production one-step implementation (D-192--D-195).
\documentclass[11pt,a4paper]{article}

\usepackage[T1]{fontenc}
\usepackage[utf8]{inputenc}
\usepackage{lmodern}
\usepackage[margin=24mm]{geometry}
\usepackage{amsmath,amssymb,bm}
\usepackage{booktabs}
\usepackage{array}
\usepackage{xcolor}
\usepackage{xurl}
\usepackage[hidelinks]{hyperref}
\usepackage{microtype}

\newcommand{\R}{\mathbb{R}}
\newcommand{\col}{\operatorname{col}}
\newcommand{\range}{\operatorname{range}}
\newcommand{\T}{\mathsf{T}}
\newcommand{\code}[1]{\texttt{#1}}
\newcolumntype{L}[1]{>{\raggedright\arraybackslash}p{#1}}

\definecolor{baselineblue}{RGB}{28,82,135}
\definecolor{annorange}{RGB}{196,92,31}
\definecolor{scopegray}{RGB}{75,75,75}

\title{From Gy\"or\"ok's Input--Output Orthogonal Augmentation\\
to One-Step State-Space OBC for the Gantry}
\author{Implementation-aligned derivation}
\date{17 September 2026}

\begin{document}
\maketitle

\begin{abstract}
This note derives the gantry's one-step OBC construction from Gy\"or\"ok's
additive input--output model.  It first translates that model to the gantry's
physical-plus-augmentation state space.  It then shows
the local sensitivity substitution required by the nonlinear LPV RK4
transition and connects each mathematical object to the production code.  The
result is a stacked, one-step orthogonality condition in normalized routed
physical-state coordinates.  It is not a multistep or closed-loop
input--output orthogonality result.
\end{abstract}

\section*{Source basis and authority}

The input--output-to-state-space derivation in Steps 1--3 follows
\path{obc-gantry-physical-state-extension.pdf}.  The choice of the local
parameter tangent, reference stack, SVD subtraction and training lifecycle in
Steps 4--8 follows \path{obc-gantry-one-step-implementation-plan.pdf}, checked
against the production implementation.  Where the implementation plan
predates a later decision, the current code and decisions D-194--D-200 are
authoritative.  In particular, the production reference set removes two
boundary samples per record, and the rollout correction uses the exact
hand-written RK4 forward sensitivity rather than a per-timestep
\code{torch.func.jvp}.

\section*{Derivation roadmap}

The construction is developed in the following order:
\begin{enumerate}
\item state Gy\"or\"ok's original input--output model and subtraction;
\item identify where the baseline and ANN add in the gantry model;
\item derive the exact linear-in-parameter state-space analogue;
\item identify the ANN rows that can overlap with the physical baseline;
\item replace Gy\"or\"ok's exact regressor by the local RK4 parameter tangent;
\item construct and stack the fixed one-step reference evaluations;
\item apply the SVD pseudoinverse subtraction;
\item insert the correction into the closed-loop rollout;
\item extend the protected space with the frozen affine-offset column; and
\item state the exact scope of the resulting orthogonality condition.
\end{enumerate}

Each step first gives the mathematical operation and then identifies the
corresponding production code.

\section{Step 1: Gy\"or\"ok's input--output construction}

Gy\"or\"ok denotes the finite-memory input--output regression vector by
\(x_k\).  To distinguish it from the gantry state later in this note, write it
as \(x_k^{\mathrm{IO}}\) without changing its meaning.  The additive predictor
is
\begin{equation}
\boxed{
\widehat y_k
=
\underbrace{\phi(x_k^{\mathrm{IO}})\theta_b}_{\text{baseline contribution}}
+
\underbrace{f_{\theta_a}^{\mathrm{ANN}}(x_k^{\mathrm{IO}})}_{\text{learning contribution}}.
}
\label{eq:gyorok-model}
\end{equation}
The important structural assumption is that the baseline is exactly linear in
its parameters:
\begin{equation}
f_{\mathrm{base}}(x_k^{\mathrm{IO}};\theta_b)
=\phi(x_k^{\mathrm{IO}})\theta_b.
\label{eq:gyorok-linear}
\end{equation}

Stacking all scalar output entries in the training data gives
\begin{equation}
\Phi=\col_{k=1}^{N}\phi(x_k^{\mathrm{IO}})\in\R^{N\times p},
\qquad
F_{\theta_a}^{\mathrm{ANN}}
=\col_{k=1}^{N}f_{\theta_a}^{\mathrm{ANN}}(x_k^{\mathrm{IO}})\in\R^N.
\end{equation}
The component of the ANN field that can be represented by the baseline is
\begin{equation}
\Phi\theta_{\mathrm{aux}},
\qquad
\theta_{\mathrm{aux}}=\Phi^\dagger F_{\theta_a}^{\mathrm{ANN}}.
\end{equation}
Gy\"or\"ok subtracts this component:
\begin{equation}
\boxed{
\widetilde F_{\theta_a}^{\mathrm{ANN}}
=F_{\theta_a}^{\mathrm{ANN}}-\Phi\theta_{\mathrm{aux}}
=(I-\Phi\Phi^\dagger)F_{\theta_a}^{\mathrm{ANN}}.
}
\label{eq:gyorok-projection}
\end{equation}
Consequently,
\begin{equation}
\boxed{\Phi^\T\widetilde F_{\theta_a}^{\mathrm{ANN}}=0.}
\label{eq:gyorok-orthogonality}
\end{equation}
At a new regression point, the stored coefficient is used as
\begin{equation}
\widetilde f_{\theta_a}^{\mathrm{ANN}}(x^{\mathrm{IO}})
=f_{\theta_a}^{\mathrm{ANN}}(x^{\mathrm{IO}})
-\phi(x^{\mathrm{IO}})\theta_{\mathrm{aux}}.
\label{eq:gyorok-prediction}
\end{equation}

In the reference implementation this lifecycle appears in
\path{orthogonal-IO-augm-main/orthogonalx_augm/src.py}, class
\code{IO\_orthogonal\_augmentation}: \(\Phi\) and its solve operator are
formed outside the objective, while \(\theta_{\mathrm{aux}}\) is recomputed
from the current ANN inside each objective evaluation.

\clearpage
\section{Step 2: Locate the addition in the gantry model}

For the gantry, the baseline and ANN do not add directly in the measured
output equation.  They add in the one-step state transition.  Define
\begin{equation}
\widehat x_k=\col(x_{b,k},x_{a,k}),
\qquad
x_{b,k}\in\R^6,\quad x_{a,k}\in\R^8,
\quad \widehat x_k\in\R^{14}.
\end{equation}
Here \(x_b=\col(q,\dot q)\) is the physical gantry state and \(x_a\) is the
learned dynamic state.
For the gantry,
\begin{equation}
q=\col(X,\Theta,Y),
\qquad
M(Y;c(v))\ddot q+C(c(v))\dot q+K(c(v))q=Pu,
\qquad Y=q_3.
\label{eq:continuous-gantry}
\end{equation}
The implemented RK4 integration of
Eq.~\eqref{eq:continuous-gantry} defines the one-sample physical transition
\begin{equation}
x_{b,k+1}=f_{\mathrm{base}}^d(v,x_{b,k},u_k).
\end{equation}

Lift the physical transition to the complete state coordinates:
\begin{equation}
F_{\mathrm{base}}^d(v,\widehat x_k,u_k)
=
\begin{bmatrix}
f_{\mathrm{base}}^d(v,x_{b,k},u_k)\\[1mm]
0
\end{bmatrix}.
\label{eq:lifted-baseline}
\end{equation}
Partition the ANN state contribution as
\begin{equation}
F_{\mathrm{aug}}(\theta_a,\widehat x_k,u_k)
=
\begin{bmatrix}
f_{\mathrm{aug}}^b(\theta_a,\widehat x_k,u_k)\\[1mm]
f_{\mathrm{aug}}^a(\theta_a,\widehat x_k,u_k)
\end{bmatrix}.
\label{eq:partitioned-ann}
\end{equation}
Before OBC, the implemented augmented model is therefore
\begin{equation}
\boxed{
\begin{aligned}
x_{b,k+1}
&=f_{\mathrm{base}}^d(v,x_{b,k},u_k)
  +f_{\mathrm{aug}}^b(\theta_a,\widehat x_k,u_k),\\
x_{a,k+1}
&=f_{\mathrm{aug}}^a(\theta_a,\widehat x_k,u_k),\\
\widehat y_k
&=h_{\mathrm{base}}(x_{b,k},u_k)
 =C_d x_{b,k}+D_d u_k.
\end{aligned}}
\label{eq:gantry-unprojected}
\end{equation}

Thus, the state-transition contribution plays the role that the direct output
contribution plays in Gy\"or\"ok's model.  The ANN has no direct output term in
Eq.~\eqref{eq:gantry-unprojected}; it affects future outputs through state
propagation.

\begin{center}
\begin{tabular}{L{0.39\textwidth}L{0.49\textwidth}}
\toprule
\textbf{Gy\"or\"ok input--output model} & \textbf{Gantry state-space adaptation}\\
\midrule
Input--output regression vector \(x_k^{\mathrm{IO}}\) & State--input point \((\widehat x_k,u_k)\)\\
Predicted output \(\widehat y_k\) & Predicted next state \(\widehat x_{k+1}\)\\
Baseline output contribution & Physical RK4 transition, lifted with zero additional-state rows\\
ANN output contribution & ANN physical-state and additional-state contributions\\
Output-space baseline span & One-step physical-transition tangent space\\
Training-regressor stack & Fixed reference-state stack\\
\bottomrule
\end{tabular}
\end{center}

\paragraph{Production code.}
Equation~\eqref{eq:gantry-unprojected} is wired in
\path{scripts/gantry/gantry_dynamic/model.py}.  The physical block receives
\((x_b,u)\) and writes the first six progressed-state rows.  The ANN receives
\((\widehat x,u)\) and its routed output is added to the progressed state.  The
physical output block receives only \((x_b,u)\).  The current production entry
uses \code{nx\_ann=8} and routes the ANN to all fourteen state rows.

The encoder supplies only the initial complete state for an actual rollout,
\begin{equation}
\widehat x_{k\mid k}
=\psi_{\theta_e}(y_{k-n_y:k-1},u_{k-n_u:k-1}).
\end{equation}
It is not used to construct the fixed OBC reference set and is not part of the
pointwise transition projection.

\clearpage
\section{Step 3: Derive the exact state-space analogue}

Before treating the nonlinear gantry, first make the translation of
Gy\"or\"ok's algebra exact.  Suppose that the physical state transition were
affine in identifiable baseline coefficients \(\vartheta_b\):
\begin{equation}
f_{\mathrm{base}}^d(\vartheta_b,x_b,u)
=f_0^d(x_b,u)+\phi_b(x_b,u)\vartheta_b.
\label{eq:hypothetical-linear-state}
\end{equation}
This is the state-equation counterpart of
Eq.~\eqref{eq:gyorok-linear}.  At evaluation point \(i\), define
\begin{equation}
\phi_{b,i}=\phi_b(x_{b,i},u_i),
\qquad
f_{\mathrm{aug},i}^b=f_{\mathrm{aug}}^b(\theta_a,\widehat x_i,u_i),
\qquad
f_{\mathrm{aug},i}^a=f_{\mathrm{aug}}^a(\theta_a,\widehat x_i,u_i).
\end{equation}
The baseline parameter directions and ANN contribution in the complete state
coordinates are
\begin{equation}
\overline\phi_{b,i}
=
\begin{bmatrix}
\phi_{b,i}\\
0
\end{bmatrix},
\qquad
f_{\mathrm{aug},i}
=
\begin{bmatrix}
f_{\mathrm{aug},i}^b\\
f_{\mathrm{aug},i}^a
\end{bmatrix}.
\label{eq:lifted-exact-directions}
\end{equation}
Their overlap is
\begin{equation}
\overline\phi_{b,i}^{\T}f_{\mathrm{aug},i}
=\phi_{b,i}^{\T}f_{\mathrm{aug},i}^b.
\label{eq:zero-block}
\end{equation}
The additional-state rows make no contribution because the baseline direction
has a zero block there.  Only \(f_{\mathrm{aug}}^b\) can overlap directly with
the baseline parameter directions.

Now stack the exact state regressors and physical ANN contributions:
\begin{equation}
\Phi_b=\col_{i=1}^{N}\phi_{b,i},
\qquad
F_{\mathrm{aug}}^b=\col_{i=1}^{N}f_{\mathrm{aug},i}^b.
\end{equation}
Applying Gy\"or\"ok's subtraction without any further change would give
\begin{equation}
\theta_{\mathrm{aux}}=\Phi_b^\dagger F_{\mathrm{aug}}^b,
\qquad
\widetilde F_{\mathrm{aug}}^b
=F_{\mathrm{aug}}^b-\Phi_b\theta_{\mathrm{aux}},
\qquad
\boxed{\Phi_b^\T\widetilde F_{\mathrm{aug}}^b=0.}
\label{eq:exact-state-obc}
\end{equation}
The additional-state update remains unchanged.  This establishes the precise
state-space version of Gy\"or\"ok's construction under the hypothetical affine
baseline representation in Eq.~\eqref{eq:hypothetical-linear-state}.

This is only a one-step coordinate statement.  A nonzero learned state can
later affect \(f_{\mathrm{aug}}^b\), the recursively propagated physical state,
and the measured output.  Equation~\eqref{eq:exact-state-obc} does not establish
multistep or input--output orthogonality.

\section{Step 4: Replace the exact regressor by the local gantry tangent}

The gantry transition is nonlinear in its physical parameters: the ten
identifiable combinations enter the mass, damping and stiffness structure and
all RK4 stages.  Hence no parameter-independent regressor \(\phi(x,u)\) gives
the complete transition globally.

The production block trains a dimensionless free-coordinate vector
\begin{equation}
v\in\R^{10}.
\end{equation}
It maps \(v\) to the identifiable physical combinations
\begin{equation}
c(v)=\col(k_{b,\mathrm{sum}},c_{g1},c_{g2},c_y,c_{b,\mathrm{sum}},m_h,
m_{\mathrm{total}},m_{\mathrm{diff}},J_{\mathrm{eff}},d).
\end{equation}
Nine coordinates use \(c_i(v)=c_{i,0}\exp(v_i)\); the signed mass difference
uses \(m_{\mathrm{diff}}(v)=m_{\mathrm{diff},0}(1+v_i)\).

At a detached epoch expansion point \(\bar v\), define
\begin{equation}
\boxed{
J_{b,k}(\bar v)
=
\left.
\frac{\partial f_{\mathrm{base}}^d(v,x_{b,k},u_k)}{\partial v}
\right|_{v=\bar v}
\in\R^{6\times10}.
}
\label{eq:local-tangent}
\end{equation}
It gives the local expansion
\begin{equation}
f_{\mathrm{base}}^d(\bar v+\Delta v,x_b,u)
\approx
f_{\mathrm{base}}^d(\bar v,x_b,u)
+J_b(\bar v,x_b,u)\Delta v.
\end{equation}
The central methodological substitution is
\begin{equation}
\boxed{
\underbrace{\phi_b(x_{b,k},u_k)}_{\substack{\text{exact state regressor}\\\text{in the affine analogue}}}
\quad\longrightarrow\quad
\underbrace{J_{b,k}(\bar v)}_{\substack{\text{local one-step}\\\text{transition sensitivity}}}.
}
\label{eq:central-substitution}
\end{equation}

This substitution preserves the projection algebra, but it does not transfer
all guarantees proved for an exact parameter-linear input--output model.

\paragraph{Production code.}
The ten free coordinates and their map are implemented by
\path{Reduced_Gantry_State_Block}.  The class is defined in
\begin{center}
\small\path{model_augmentation/fit_systems/blocks.py}.
\end{center}
Its
\code{transition\_from\_free(v,z)} evaluates the complete normalized RK4
transition.  The reference basis is differentiated with
\code{torch.func.jacfwd} by \code{build\_stacked\_sensitivity} in
\path{model_augmentation/fit_systems/obc.py}.

\clearpage
\section{Step 5: Construct and stack the fixed reference evaluations}

The construction uses every valid sample from all fourteen training records.
From measured stage positions,
\begin{equation}
q_k^{\mathrm{ref}}=P^{-\T}y_k^{\mathrm{data}}.
\end{equation}
The unmeasured velocity is estimated inside each record using
\begin{equation}
\widehat{\dot q}_k^{\mathrm{FD}}
=\frac{-q_{k+2}^{\mathrm{ref}}+8q_{k+1}^{\mathrm{ref}}
-8q_{k-1}^{\mathrm{ref}}+q_{k-2}^{\mathrm{ref}}}{12T_s}.
\label{eq:fd}
\end{equation}
Two samples are removed at each end of every record, so the stencil never
crosses a record boundary.  Define
\begin{equation}
x_{b,k}^{\mathrm{ref}}
=\col(q_k^{\mathrm{ref}},\widehat{\dot q}_k^{\mathrm{FD}}),
\qquad
x_{a,k}^{\mathrm{ref}}=0,
\qquad
u_k^{\mathrm{ref}}=u_k^{\mathrm{data}}.
\end{equation}
These tuples are fixed one-step evaluation points, not a simulated trajectory.
The controller is not differentiated: \(u_k^{\mathrm{ref}}\) is held fixed.

Let \(R_b\in\R^{n_r\times6}\) select the physical rows written by the ANN.
At every reference point,
\begin{equation}
\mathcal J_k
=R_bJ_{b,k}^{\mathrm n}(\bar v)
\in\R^{n_r\times10},
\end{equation}
and
\begin{equation}
g_k(\theta_a)
=R_b f_{\mathrm{aug}}^{b,\mathrm n}
(\theta_a,x_{b,k}^{\mathrm{ref}},0,u_k^{\mathrm{ref}})
\in\R^{n_r}.
\end{equation}
The superscript \({\mathrm n}\) emphasizes that both quantities are already
in the model's normalized state coordinates.

Stack the samples in sample-major order:
\begin{equation}
\boxed{
J
=\begin{bmatrix}\mathcal J_1\\ \vdots\\ \mathcal J_N\end{bmatrix}
\in\R^{Nn_r\times10},
\qquad
F_{\mathrm{aug}}(\theta_a)
=\begin{bmatrix}g_1\\ \vdots\\ g_N\end{bmatrix}
\in\R^{Nn_r}.
}
\label{eq:gantry-stacks}
\end{equation}
For the production reference set, \(N=671{,}944\) and all six physical rows
are routed, so
\begin{equation}
J\in\R^{4{,}031{,}664\times10}.
\end{equation}

\paragraph{Production code.}
Reference construction is in
\path{scripts/gantry/gantry_dynamic/obc_gantry.py}:
\code{build\_reference\_set}.  The basis stack is created by
\code{make\_basis\_builder}; the matching ANN stack is created by
\code{make\_reference\_field}.  Both use the same physical-row order and
sample-major flattening.

\clearpage
\section{Step 6: Apply Gy\"or\"ok's subtraction to the gantry stack}

Factor the stacked tangent directly with an economy SVD:
\begin{equation}
J=USV^\T.
\end{equation}
The implementation does not invert the Gram matrix.  It computes
\begin{equation}
\boxed{
\theta_{\mathrm{aux}}(\theta_a)
=J^\dagger F_{\mathrm{aug}}(\theta_a)
=VS^{-1}U^\T F_{\mathrm{aug}}(\theta_a)
\in\R^{10}.
}
\label{eq:gantry-aux}
\end{equation}
The corrected stacked physical ANN field is
\begin{equation}
\boxed{
\widetilde F_{\mathrm{aug}}
=F_{\mathrm{aug}}-J\theta_{\mathrm{aux}}.
}
\label{eq:gantry-corrected-stack}
\end{equation}
Therefore,
\begin{equation}
\boxed{J^\T\widetilde F_{\mathrm{aug}}=0}
\label{eq:gantry-orthogonality}
\end{equation}
to numerical precision.

The correspondence with Gy\"or\"ok is now explicit:
\begin{equation}
\begin{aligned}
\Phi
&\longrightarrow J,\\
\Phi^\dagger F_{\mathrm{ANN}}
&\longrightarrow J^\dagger F_{\mathrm{aug}},\\
F_{\mathrm{ANN}}-\Phi\theta_{\mathrm{aux}}
&\longrightarrow F_{\mathrm{aug}}-J\theta_{\mathrm{aux}}.
\end{aligned}
\end{equation}

\paragraph{Production code.}
\code{OBCBasis.from\_matrix} in
\path{model_augmentation/fit_systems/obc.py} performs the float64 economy
SVD.  \code{OBCBasis.coefficient} evaluates
\(VS^{-1}U^\T F_{\mathrm{aug}}\) without detaching the ANN field, so gradients
can pass through \(\theta_{\mathrm{aux}}\) to the ANN.

\section{Step 7: Respect the training lifecycle}

At the beginning of epoch \(r\), the code takes a detached copy of the current
physical free coordinate,
\begin{equation}
\bar v^{(r)}=v^{(r)},
\end{equation}
rebuilds \(J^{(r)}\), and holds its SVD factors fixed for that epoch.

At every objective evaluation, it evaluates the current ANN on the fixed
reference set and recomputes
\begin{equation}
\theta_{\mathrm{aux}}(\theta_a)
=(J^{(r)})^\dagger F_{\mathrm{aug}}(\theta_a).
\end{equation}
The coefficient is computed once per objective, remains differentiable, and is
reused at every timestep of that objective's rollout.  It is never reused
across optimizer updates.

This lifecycle is implemented by \code{OBCLifecycle} in
\path{model_augmentation/fit_systems/obc.py} and invoked by
\code{SSE\_Interconnect\_Composed.loss} in
\path{model_augmentation/fit_systems/interconnect.py}.

\clearpage
\section{Step 8: Insert the correction into the rollout model}

At an actual rollout point \((\widehat x_k,u_k)\), evaluate the local
directional correction using the epoch expansion point:
\begin{equation}
c_k
=J_b^{\mathrm n}(\bar v^{(r)},x_{b,k},u_k)
\theta_{\mathrm{aux}}.
\end{equation}
The OBC model executed by the rollout is
\begin{equation}
\boxed{
\begin{aligned}
x_{b,k+1}
&=f_{\mathrm{base}}^{d,\mathrm n}(v,x_{b,k},u_k)
 +f_{\mathrm{aug}}^{b,\mathrm n}(\theta_a,\widehat x_k,u_k)
 -J_b^{\mathrm n}(\bar v^{(r)},x_{b,k},u_k)\theta_{\mathrm{aux}},\\
x_{a,k+1}
&=f_{\mathrm{aug}}^{a,\mathrm n}(\theta_a,\widehat x_k,u_k),\\
\widehat y_k
&=h_{\mathrm{base}}(x_{b,k},u_k).
\end{aligned}}
\label{eq:corrected-rollout}
\end{equation}

The controller remains active and determines the inputs encountered during the
closed-loop rollout.  Its derivative is not part of \(J\) or
Eq.~\eqref{eq:gantry-orthogonality}.

\paragraph{Current derivative implementation.}
The mathematical definition of the correction is the directional derivative
\(J_b\theta_{\mathrm{aux}}\).  Production code evaluates it with the exact
hand-written forward sensitivity through the RK4 stages.  The
parameter-to-matrix tangent is formed once per objective and threaded into the
rollout.  \code{torch.func.jvp} is retained as a test oracle, not used in the
compiled per-timestep production path.  The relevant implementations are
\code{GantryOBCCorrection.tangent\_pair} and
\code{GantryOBCCorrection.forward} in
\path{scripts/gantry/gantry_dynamic/obc_gantry.py}, together with
\code{\_rk4\_with\_tangent} in
\path{model_augmentation/fit_systems/blocks.py}.

The subtraction itself is applied only to the progressed-state signal in
\code{Interconnect.forward}; the output map is untouched.

\clearpage
\section{Step 9: Add the affine-offset comparison arm}

The primary implementation protects the ten-column tangent \(J\). Decision
D-200 adds a separate arm that also protects the frozen baseline transition:
\begin{equation}
\boxed{
B=\begin{bmatrix}J&c\end{bmatrix},\qquad
c=\col_i R_b f_{\mathrm{base}}^{d,\mathrm n}
   (\bar v,z_i^{\mathrm{ref}})\in\R^{N_rn_r}.
}
\label{eq:affine-basis}
\end{equation}
This is the same column space as the centered affine form
\([J\mid\Gamma]\), because \(\Gamma=c-J\bar v\). Hence
\begin{equation}
\range([J\mid c])=\range([J\mid\Gamma]).
\end{equation}
The code uses \(c\), since the RK4 forward-sensitivity calculation already
returns the frozen primal transition.

After column equilibration and a float64 SVD of \(B\), compute
\begin{equation}
\begin{bmatrix}a\\ \alpha\end{bmatrix}=B^\dagger F_{\mathrm{aug}},
\qquad
\widetilde F_{\mathrm{aug}}=F_{\mathrm{aug}}-Ja-\alpha c,
\qquad
B^\T\widetilde F_{\mathrm{aug}}=0.
\end{equation}
At an arbitrary rollout state, the implemented correction is
\begin{equation}
\boxed{
J_b(\bar v,x_{b,k},u_k)a
+\alpha f_{\mathrm{base}}^{d,\mathrm n}(\bar v,x_{b,k},u_k).
}
\label{eq:affine-runtime}
\end{equation}
Both terms use the epoch's frozen \(\bar v\). Using the live physical
parameters in the second term would evaluate a different column from the one
used to fit \(\alpha\), so the reference-stack identity would no longer match
the rollout implementation. The optimized RK4 routine returns the primal and
directional derivative together; Eq.~\eqref{eq:affine-runtime} therefore adds
no second RK4 call per timestep.

\paragraph{Production code.}
Set \code{OBC\_ARM=obc\_affine}. The imported configuration records
\code{obc=True} and \code{obc\_space='affine'}. The existing
\code{OBC\_ARM=obc} remains the unscaled ten-column tangent arm.

\clearpage
\section{Step 10: State precisely what has been established}

The implemented condition is
\begin{quote}
The current ANN's stacked one-step contribution in normalized routed physical
state coordinates is orthogonal to the selected local protected space: the
ten-dimensional parameter tangent for \code{obc}, or the eleven-column tangent
plus frozen transition for \code{obc\_affine}. Both are evaluated at the epoch
expansion point on the fixed reference slice \(x_a=0\).
\end{quote}

More compactly,
\begin{equation}
\boxed{
J(\bar v;\mathcal Z_{\mathrm{ref}})^\T
\left[
F_{\mathrm{aug}}(\theta_a;\mathcal Z_{\mathrm{ref}})
-J(\bar v;\mathcal Z_{\mathrm{ref}})\theta_{\mathrm{aux}}
\right]=0.
}
\end{equation}

It does \emph{not} establish:
\begin{itemize}
\item pointwise orthogonality at every individual reference point;
\item global orthogonality over the nonlinear physical parameter space;
\item orthogonality away from the reference slice \(x_a=0\);
\item orthogonality of recursively propagated or multistep state sensitivities;
\item closed-loop input--output Jacobian orthogonality;
\item stability, covariance preservation, consistency, or unique physical
      parameter recovery under Gy\"or\"ok's theorems.
\end{itemize}

The construction intentionally follows the requested scope: retain the normal
physical-plus-augmentation state, introduce no delayed input/output state, and
first test the regular one-step orthogonalization before considering a
different level of orthogonality.

\section{One-line derivation chain}

\begin{equation}
\boxed{
\begin{gathered}
\text{Gy\"or\"ok: }\Phi
\quad\longrightarrow\quad
\text{gantry: }J
=\col_k R_b
\left.\frac{\partial f_{\mathrm{base}}^{d,\mathrm n}}
{\partial v}\right|_{\bar v,z_k^{\mathrm{ref}}}
\\[1mm]
J=USV^\T,
\qquad
\theta_{\mathrm{aux}}=VS^{-1}U^\T F_{\mathrm{aug}},
\\[1mm]
\widetilde f_{\mathrm{aug},k}^b
=f_{\mathrm{aug},k}^b
-J_b(\bar v,x_{b,k},u_k)\theta_{\mathrm{aux}},
\\[1mm]
J^\T\widetilde F_{\mathrm{aug}}=0
\quad\text{on the fixed one-step reference stack.}
\end{gathered}}
\end{equation}

\end{document}
